\documentclass[11pt]{article}

\usepackage[T1]{fontenc}
\usepackage{lmodern}
\usepackage{microtype}
\usepackage[a4paper,margin=25mm,headheight=14pt]{geometry}
\usepackage{booktabs,longtable,array,tabularx}
\usepackage{enumitem}
\usepackage{xcolor}
\usepackage{hyperref}
\usepackage{xurl}
\usepackage{fancyhdr}
\usepackage{listings}

\setlength{\emergencystretch}{3em}
\newcolumntype{L}[1]{>{\raggedright\arraybackslash}p{#1}}

\ifdefined\pdfinfoomitdate
  \pdfinfoomitdate=1
\fi
\ifdefined\pdftrailerid
  \pdftrailerid{}
\fi
\ifdefined\pdfsuppressptexinfo
  \pdfsuppressptexinfo=15
\fi

\hypersetup{
  colorlinks=true,
  linkcolor=blue!55!black,
  urlcolor=blue!55!black,
  pdftitle={Formal Tool Adoption Audit},
  pdfsubject={Adoption decision record for formal and certified-numerical tools proposed for pid-rs},
  pdfkeywords={pid-rs, Verus, Kani, Aeneas, MPFR, Rug, Rocq, Interval, assurance},
  pdfauthor={pid-rs project analysis}
}

\lstset{
  basicstyle=\ttfamily\small,
  breaklines=true,
  breakatwhitespace=false,
  columns=fullflexible,
  keepspaces=true,
  frame=single,
  framerule=0.3pt,
  xleftmargin=0.5em,
  xrightmargin=0.5em,
  aboveskip=0.7em,
  belowskip=0.7em,
  showstringspaces=false
}

\pagestyle{fancy}
\fancyhf{}
\fancyhead[L]{pid-rs formal tool adoption audit}
\fancyhead[R]{decision record}
\fancyfoot[C]{\thepage}

\title{\textbf{Formal Tool Adoption Audit}\\
\large Verus, Kani, Aeneas, MPFR through Rug, and Rocq Interval}
\author{pid-rs project analysis}
\date{24 July 2026}

\begin{document}
\maketitle

\begin{abstract}
This document records an adoption decision for five assurance layers proposed for
\texttt{pid-rs}.  It separates upstream facts from local observations.  It records exact release
pins, licenses, toolchains, commands, blockers, trust boundaries, negative controls, and permitted
claim text.  The selected order is: a Rug and MPFR categorical interval oracle, bounded Kani
harnesses, a Verus proof pilot for the same pure categorical kernel, an optional Rocq Interval
certificate lane, and a deferred Aeneas translation experiment.  This document is not evidence
that any listed tool verified \texttt{pid-rs}.  It adds no dependency, CI gate, proof result,
certified result, or release claim.
\end{abstract}

\begingroup
\small
\setlength{\parskip}{0pt}
\tableofcontents
\endgroup

\section{Record status and purpose}

\begin{quote}
\textbf{Record status:} Adoption decision record, dated 24 July 2026.

\textbf{Evidence status:} No tool evaluated in this record has verified \texttt{pid-rs}.  This
record does not report a proof run, model-checking result, certified numerical result, or
independent certificate check.  It selects bounded pilot work and states the conditions that must
apply before a result can become repository evidence.
\end{quote}

The audit covers:
\begin{enumerate}[leftmargin=*]
  \item Verus;
  \item Kani;
  \item Aeneas with Charon;
  \item MPFR through Rug; and
  \item Rocq with the Interval package.
\end{enumerate}

The cut-off date is 24 July 2026.  Upstream state can change after that date.  A pilot must repeat
the upstream release check before installation.

The decision is:
\begin{enumerate}[leftmargin=*]
  \item Pilot a Rug and MPFR categorical reference backend first.
  \item Pilot Kani on bounded categorical integer and lattice obligations.
  \item Pilot Verus on the same small, pure categorical kernel.
  \item Evaluate Rocq Interval after the executable interval format is stable.
  \item Defer Aeneas production adoption.  Permit one isolated translation experiment.
\end{enumerate}

No one tool replaces Lean.  The tools cover different failure classes.  Evidence from one layer
does not imply evidence from another layer.

\section{Scope, exclusions, and terms}

This record covers tool fit, versions, licenses, toolchains, installation pins, reproducible
commands, trust boundaries, and permitted claim language.

This record does not:
\begin{itemize}[leftmargin=*]
  \item add a dependency or change CI;
  \item change a public API, estimator, method catalog, or scientific claim;
  \item prove Rust-to-specification refinement;
  \item prove floating-point correctness or statistical calibration;
  \item prove a theorem about continuous PID;
  \item approve a license for distribution; or
  \item claim that the current local installation is a release environment.
\end{itemize}

License identifiers come from upstream metadata.  They are not legal advice.  A distribution
review is mandatory before code that links to an LGPL or CeCILL-C component is released.

\paragraph{Upstream fact.}
A fact in an official project repository, release page, manual, or package record.

\paragraph{Local fact.}
An observation from the workstation on 24 July 2026.

\paragraph{Trust boundary.}
Code, data, assumptions, or tools that a proof result does not verify.  This document also uses
the common term trusted computing base, or TCB.

\paragraph{Pilot.}
A bounded evaluation.  It does not mean that a tool is part of the build, CI, or release process.

\paragraph{Certified backend.}
A new backend that returns an outward interval under a stated rounding contract.  It does not mean
that an existing binary64 result is certified.

\section{Current repository and local state}

\subsection{Repository pins}

\begin{center}
\begin{tabular}{L{0.34\textwidth}L{0.22\textwidth}L{0.34\textwidth}}
\toprule
Item & Repository pin & Source \\
\midrule
Minimum supported Rust version & 1.89 & \path{Cargo.toml} \\
Lean & 4.32.0 & \path{audit/formal/lean/lean-toolchain} \\
mathlib & \path{v4.32.0} & \path{audit/formal/lean/lakefile.toml} \\
Z3 for the existing algebra checker & 4.16.0, 64 bit &
\path{scripts/check-z3-pid2-algebra.py} \\
\bottomrule
\end{tabular}
\end{center}

These pins do not apply automatically to a new verifier.  Each verifier can have a different
compiler and solver pin.

\subsection{Local observations}

The following commands were run from the repository worktree:

\begin{lstlisting}
rustc --version
cargo --version
(cd audit/formal/lean && lake env lean --version)
z3 --version
pkg-config --modversion mpfr
pkg-config --modversion gmp
command -v <tool>
\end{lstlisting}

\begin{longtable}{L{0.34\textwidth}L{0.56\textwidth}}
\toprule
Local item & Observation \\
\midrule
\endfirsthead
\toprule
Local item & Observation \\
\midrule
\endhead
\path{rustc} & \path{rustc 1.96.0 (ac68faa20 2026-05-25)} \\
\path{cargo} & \path{cargo 1.96.0 (30a34c682 2026-05-25)} \\
Lean & 4.32.0, \path{arm64-apple-darwin24.6.0}, commit
\path{8c9756b28d64dab099da31a4c09229a9e6a2ef35} \\
Z3 & 4.16.0, 64 bit \\
MPFR found by \path{pkg-config} & 4.2.2 \\
GMP found by \path{pkg-config} & 6.3.0 \\
FLINT found by \path{pkg-config} & Not found \\
\path{rug} or \path{gmp-mpfr-sys} in \path{Cargo.lock} & Not present \\
\bottomrule
\end{longtable}

The commands \path{verus}, \path{cargo-verus}, \path{kani}, \path{cargo-kani}, \path{charon},
\path{aeneas}, \path{opam}, \path{rocq}, \path{rocqchk}, \path{coqc}, \path{coqchk}, and \path{nix}
were not found.

This local list is an environment observation only.  It is not an installation attestation.
Homebrew MPFR availability does not establish a Rust binding, directed-rounding wrapper, or
certificate path.

\section{Decision matrix}

\small
\begin{longtable}{
  L{0.12\textwidth}
  L{0.10\textwidth}
  L{0.19\textwidth}
  L{0.18\textwidth}
  L{0.19\textwidth}}
\toprule
Layer & Decision & First permitted scope & Main reason & Main blocker \\
\midrule
\endfirsthead
\toprule
Layer & Decision & First permitted scope & Main reason & Main blocker \\
\midrule
\endhead
Rug with MPFR & Pilot 1 & Offline or optional categorical exact-count interval oracle &
Direct control of log rounding and atom cancellation &
Native C and FFI trust, wrapper review, and LGPL obligations \\
Kani & Pilot 2 & Bounded integer, mask, event, lattice, overflow, and rejection harnesses &
Exhaustive bit-precise checks inside declared bounds &
Bounds do not generalize; concurrency and transcendental accuracy are out of scope \\
Verus & Pilot 3 & Small pure categorical kernel with explicit contracts &
Closest fit for unbounded functional contracts over Rust-like code &
Rust 1.96 versus the repository Rust 1.89 minimum; unsupported code and assumptions expand the
TCB \\
Rocq Interval & Pilot 4* &
Independent checks of fixed emitted interval inequalities &
Small kernel checker and analytic interval tactics give a second proof route &
Separate ecosystem; the certificate generator and runtime binding remain trusted \\
Aeneas and Charon & Defer; one experiment & Translation of the same small pure kernel to Lean &
Can connect accepted safe Rust to a theorem-prover model &
No stable release, Lean 4.31 mismatch, supported-subset limits, and handwritten external models \\
\bottomrule
\end{longtable}
\normalsize

* Pilot 4 is conditional on a stable, versioned certificate schema.

The order is deliberate.  The first pilot gives a numerical oracle that the other pilots can use.
Kani and Verus then target the same integer semantics with different proof models.  Rocq must not
precede a stable certificate schema.  Aeneas must not become a release dependency during this
evaluation.

\section{Verus adoption record}

\subsection{Upstream pin and license}

The selected non-prerelease cut is
\href{https://github.com/verus-lang/verus/releases/tag/release/0.2026.07.18.3a4d30b}
{\path{release/0.2026.07.18.3a4d30b}}.  GitHub records its publication on 20 July 2026.

The release pins:
\begin{itemize}[leftmargin=*]
  \item Rust 1.96.0;
  \item \path{vstd} version \path{0.0.0-2026-07-12-0122}.
\end{itemize}

The Z3 fetch helper in the tagged source pins Z3 4.12.5.

The license is MIT.  At the audit date, the release page had assets for arm64 macOS, x86-64
macOS, x86-64 Linux, and x86-64 Windows.  This list is not a general platform support theorem.

Primary pin sources are the
\href{https://raw.githubusercontent.com/verus-lang/verus/release/0.2026.07.18.3a4d30b/rust-toolchain.toml}
{release toolchain},
\href{https://raw.githubusercontent.com/verus-lang/verus/release/0.2026.07.18.3a4d30b/source/vstd/Cargo.toml}
{\path{vstd} package metadata},
\href{https://raw.githubusercontent.com/verus-lang/verus/release/0.2026.07.18.3a4d30b/.github/workflows/get-z3.sh}
{Z3 fetch helper}, and
\href{https://raw.githubusercontent.com/verus-lang/verus/release/0.2026.07.18.3a4d30b/LICENSE}
{MIT license}.

\subsection{Pin and installation rule}

Do not use a rolling prerelease tag.  Do not mix \path{vstd}, the Verus executable, and solver
files from different releases.

For a pilot:
\begin{enumerate}[leftmargin=*]
  \item Select the platform asset from the immutable release page.
  \item Compute and record its SHA-256.
  \item Extract it into a version-specific directory.
  \item Use \path{cargo-verus} from that release.
  \item Use the solver in the release bundle, or the exact upstream solver pin.
  \item Record all tool versions, asset digest, host, and target.
\end{enumerate}

The Cargo metadata and commands are:

\begin{lstlisting}
[package.metadata.verus]
verify = true

cargo verus verify --workspace --locked
cargo verus build --release --locked
\end{lstlisting}

Verus also supports \path{--offline}.  A release evidence job must use a lockfile and an archived
dependency set before it claims offline reproducibility.  A focused verification run is a
development aid.  It is not a substitute for the full verification command.

\subsection{First obligations}

The pilot can cover:
\begin{itemize}[leftmargin=*]
  \item multiset preservation by sorting and run-length encoding;
  \item histogram count totals;
  \item equality of event counts and event predicates;
  \item denominator positivity for supported keyed rows;
  \item canonical source-mask and antichain order;
  \item M{\"o}bius reconstruction; and
  \item specialized and general categorical-kernel equivalence.
\end{itemize}

The pilot must not start with logarithms, binary64 error, the complete report API, parallel code,
cancellation plumbing, allocators, Python bindings, or continuous estimators.

\subsection{Trust boundary, blockers, and negative controls}

Verus does not support all Rust code and libraries.  The Verus verifier, Rust compiler, LLVM,
solver, and platform remain in the trust boundary.

The features \path{assume}, \path{external_body}, \path{external_fn_specification}, and
\path{external} add assumptions or trusted specifications.  The pilot must create an assumption
inventory.  A result with an unreviewed assumption is not accepted.

Verus uses Rust 1.96.  The repository supports Rust 1.89.  A successful Verus run does not prove
that the erased program builds at Rust 1.89.  The normal Rust 1.89 gates remain separate.

Required negative controls must:
\begin{itemize}[leftmargin=*]
  \item mutate an event union to an intersection;
  \item mutate a source mask;
  \item remove a M{\"o}bius predecessor;
  \item add an unapproved assumption; and
  \item run the erased source under Rust 1.89.
\end{itemize}
The first three mutations must fail a contract.  The assumption mutation must fail policy.  The
MSRV run must detect incompatible erased source.

\subsection{Permitted claim}

\begin{quote}
Verus proved the stated contracts for the opted-in functions under the recorded assumptions and
pinned Verus, solver, and toolchain.  It did not verify rustc or LLVM, unchecked specifications or
external bodies, callers outside those contracts, or floating-point transcendental correctness.
\end{quote}

Do not state that Verus verified \texttt{pid-rs}, the Rust binary, numerical correctness, or
statistical correctness.

\section{Kani adoption record}

\subsection{Upstream pin and license}

The selected stable release is
\href{https://github.com/model-checking/kani/releases/tag/kani-0.67.0}{\path{kani-0.67.0}},
published on 16 January 2026.  It pins Rust \path{nightly-2025-11-21} and CBMC 6.8.0.
Kani is dual licensed under MIT or Apache-2.0.

Primary pin sources are:
\begin{itemize}[leftmargin=*]
  \item
  \href{https://raw.githubusercontent.com/model-checking/kani/kani-0.67.0/rust-toolchain.toml}
  {Rust toolchain};
  \item
  \href{https://raw.githubusercontent.com/model-checking/kani/kani-0.67.0/kani-dependencies}
  {dependency pins};
  \item
  \href{https://raw.githubusercontent.com/model-checking/kani/kani-0.67.0/LICENSE-MIT}
  {MIT license}; and
  \item
  \href{https://raw.githubusercontent.com/model-checking/kani/kani-0.67.0/LICENSE-APACHE}
  {Apache-2.0 license}.
\end{itemize}

The official easy-install matrix lists x86-64 Linux with GNU libc, x86-64 macOS, and arm64 macOS.
The installation guide states Rust 1.58 or newer through \path{rustup}.

\subsection{Pin and installation rule}

Use a version-specific \path{KANI_HOME}:

\begin{lstlisting}
export KANI_HOME=/absolute/tool-cache/kani/0.67.0
cargo install --locked kani-verifier --version 0.67.0
cargo kani setup
cargo kani --version
\end{lstlisting}

Archive or hash the resulting tool directory.  Record the package checksum, all versions, host,
target, harness, input bounds, unwind bound, and unwinding assertion result.

\begin{lstlisting}
cargo kani --harness <HARNESS_NAME> --unwind <BOUND>
\end{lstlisting}

\subsection{First obligations}

Use Kani for bounded checks of no panic, no out-of-bounds index, no checked integer overflow,
resource-estimate arithmetic, mask and lattice indexing, malformed-shape rejection, cancellation
paths, specialized and general kernel equivalence, and exact reconstruction identities.

Each harness must declare all shape and count bounds.  It must retain an unwinding assertion.

\subsection{Trust boundary, blockers, and negative controls}

A result for bound \(N\) does not imply a result for \(N+1\).  Kani does not model supported
concurrent execution as a concurrent proof.  Its floating-point support includes partial or
overapproximated intrinsics.  It is not a precision proof route for logarithms.

The Kani compiler, CBMC, solver backend, harness assumptions, and environment models remain in the
trust boundary.

Required negative controls must:
\begin{itemize}[leftmargin=*]
  \item reduce an unwind bound and trigger its assertion;
  \item introduce an index error and obtain a counterexample;
  \item remove an overflow check and obtain a counterexample;
  \item change a specialized mask and obtain an equivalence counterexample; and
  \item label one case outside the proved bound as not covered.
\end{itemize}

\subsection{Permitted claim}

\begin{quote}
Kani exhaustively checked this harness within the declared input and unwind bounds using Kani
0.67.0 and CBMC 6.8.0.  It does not establish correctness outside those bounds, concurrent
execution semantics, or numerical precision of transcendental functions.
\end{quote}

Do not state that Kani proved the algorithm for all inputs, verified parallel execution, or proved
the accuracy of \path{ln} or digamma.

\section{Aeneas and Charon adoption record}

\subsection{Upstream pin and license}

Aeneas had no stable release at the audit date.  Its release page contained nightly prereleases.
This record therefore uses:
\begin{itemize}[leftmargin=*]
  \item Aeneas commit
  \href{https://github.com/AeneasVerif/aeneas/commit/6e167c9b63a4dafd66d0e0edd9d94669f957ff7b}
  {\path{6e167c9b63a4dafd66d0e0edd9d94669f957ff7b}};
  \item Charon commit \path{527ea8e3b5dcb52edd6aef0f7bc34cc09c11dd59};
  \item Charon Rust nightly \path{nightly-2026-06-01};
  \item Lean 4.31.0;
  \item mathlib \path{v4.31.0}; and
  \item OCaml switch example 5.3.0.
\end{itemize}

Aeneas and Charon are licensed under Apache-2.0.

Primary pin sources are:
\begin{itemize}[leftmargin=*]
  \item \href{https://github.com/AeneasVerif/aeneas}{repository};
  \item \href{https://github.com/AeneasVerif/aeneas/releases}{release page};
  \item
  \href{https://raw.githubusercontent.com/AeneasVerif/aeneas/6e167c9b63a4dafd66d0e0edd9d94669f957ff7b/charon-pin}
  {Charon pin};
  \item
  \href{https://raw.githubusercontent.com/AeneasVerif/charon/527ea8e3b5dcb52edd6aef0f7bc34cc09c11dd59/charon/rust-toolchain}
  {Charon toolchain};
  \item
  \href{https://raw.githubusercontent.com/AeneasVerif/charon/527ea8e3b5dcb52edd6aef0f7bc34cc09c11dd59/LICENSE.md}
  {Charon Apache-2.0 license};
  \item
  \href{https://raw.githubusercontent.com/AeneasVerif/aeneas/6e167c9b63a4dafd66d0e0edd9d94669f957ff7b/backends/lean/lean-toolchain}
  {Lean toolchain};
  \item
  \href{https://raw.githubusercontent.com/AeneasVerif/aeneas/6e167c9b63a4dafd66d0e0edd9d94669f957ff7b/backends/lean/lakefile.lean}
  {mathlib pin}; and
  \item
  \href{https://raw.githubusercontent.com/AeneasVerif/aeneas/6e167c9b63a4dafd66d0e0edd9d94669f957ff7b/LICENSE.md}
  {license}.
\end{itemize}

\subsection{Reproducible experiment rule}

Use a detached full commit.  Do not use a moving nightly name as evidence identity.

\begin{lstlisting}
git clone https://github.com/AeneasVerif/aeneas.git
git -C aeneas checkout --detach 6e167c9b63a4dafd66d0e0edd9d94669f957ff7b
cd aeneas
opam switch create 5.3.0
eval "$(opam env)"
make setup-charon
make
\end{lstlisting}

\path{make setup-charon} must resolve the exact Charon commit.  Export the complete OPAM switch.
Retain \path{flake.lock} if Nix is used.

\begin{lstlisting}
charon cargo --preset=aeneas
./bin/aeneas -backend lean [OPTIONS] <LLBC_FILE>
\end{lstlisting}

Retain both commits, Charon nightly, OPAM export, Lean and mathlib pins, translation flags, source
digest, LLBC digest, generated Lean digest, and Lean checker output.

\subsection{Experiment scope and blockers}

Translate only the same small pure kernel selected for Verus.  Compare the generated statement
with the handwritten Lean specification.  The experiment must identify rejected constructs,
handwritten external models, theorem correspondence, and toolchain compatibility.

Aeneas supports a subset of safe Rust.  Unsafe code and concurrency are out of scope.  Some loop
control patterns and library operations are unsupported.  Handwritten models become trusted
inputs.

The backend pins Lean and mathlib 4.31.0.  The repository pins 4.32.0.  This mismatch blocks direct
integration.

Charon translation, Aeneas translation, handwritten models, theorem correspondence, Lean, and the
compiler path remain in the trust boundary.

Required negative controls must submit an unsupported construct, mutate an event predicate, remove
a handwritten model, alter the Charon checkout, and compare Lean 4.31.0 with 4.32.0.  Rejection
must be explicit.  A compatibility experiment must not be reported as toolchain equivalence.

\subsection{Permitted claim}

\begin{quote}
Aeneas translated the accepted safe-Rust subset at the pinned Aeneas and Charon commits, and Lean
checked the generated model's stated theorem.  This does not establish translation support for
all Rust, correctness of handwritten external models, unsafe or concurrent code, or binary
equivalence unless those links are separately proved.
\end{quote}

Do not state that Aeneas verified the Rust implementation, that generated Lean proves the compiled
binary, or that the translation covers unsupported library code.

\section{Rug and MPFR adoption record}

\subsection{Upstream pin and license}

The selected pins are:
\begin{itemize}[leftmargin=*]
  \item Rug 1.30.0, published on 10 July 2026;
  \item \path{gmp-mpfr-sys} 1.7.1;
  \item MPFR 4.2.2;
  \item GMP 6.3.0; and
  \item MPC 1.4.1 only if a selected feature requires it.
\end{itemize}

Rug requires Rust 1.85 or newer.  \path{gmp-mpfr-sys} requires Rust 1.71 or newer.  Both fit the
repository Rust 1.89 minimum.

Rug, \path{gmp-mpfr-sys}, and MPFR use LGPL-3.0-or-later licensing in the cited package records.
GMP 6.3.0 is available under LGPLv3 or GPLv2.  A distribution review is mandatory.

Primary sources are:
\begin{itemize}[leftmargin=*]
  \item \href{https://www.mpfr.org/mpfr-current/}{MPFR current release};
  \item \href{https://www.mpfr.org/mpfr-current/mpfr.html}{MPFR manual};
  \item \href{https://docs.rs/rug/1.30.0/rug/}{Rug 1.30.0 metadata};
  \item \href{https://docs.rs/rug/1.30.0/rug/float/enum.Round.html}{Rug rounding modes};
  \item \href{https://docs.rs/gmp-mpfr-sys/1.7.1/gmp_mpfr_sys/}
  {\path{gmp-mpfr-sys} metadata}; and
  \item \href{https://gmplib.org/manual/Copying}{GMP copying conditions}.
\end{itemize}

MPFR states that most functions are correctly rounded as if computed at infinite precision and
then rounded in the requested direction.  \path{mpfr_log} computes the natural logarithm.
\path{MPFR_RNDD} rounds toward negative infinity.  \path{MPFR_RNDU} rounds toward positive
infinity.  Experimental faithful rounding \path{MPFR_RNDF} must not be used for a certificate.

\subsection{Future dependency pin}

This record does not add a dependency.  A future pilot must use:

\begin{lstlisting}
[dependencies.rug]
version = "=1.30.0"
default-features = false
features = ["integer", "float", "std"]
\end{lstlisting}

Do not enable experimental \path{use-system-libs} for reproducible evidence.  Use bundled GMP and
MPFR sources.  Record Cargo checksums and native archive digests.  Use an optional offline oracle
or a separate small crate.  Do not silently replace the stable binary64 API.

\subsection{First numerical obligations}

For categorical SxPID:
\begin{enumerate}[leftmargin=*]
  \item Keep histogram and event counts as exact integers.
  \item Form exact rational count ratios.
  \item Evaluate each logarithm toward both negative and positive infinity.
  \item Apply outward rounding to later arithmetic.
  \item Apply the concrete M{\"o}bius transform to intervals.
  \item Increase precision until a declared width or stability rule applies.
  \item Return \path{Unresolved} if an interval contains zero.
  \item Return an ordering ambiguity if \(\mathrm{I}_{\min}\) candidate intervals overlap.
  \item Retain exact input, expression-schema, precision-policy, and output digests.
\end{enumerate}

The schema must distinguish certified positive, certified negative, certified exact zero,
unresolved sign, unresolved minimizer order, and resource or precision limit.  A midpoint is not a
certificate.  A narrow interval is not proof of statistical validity.

This first backend does not cover continuous KSG.  KSG needs separate proofs for neighbor
selection, represented distances, strict-radius semantics, estimator assumptions, and statistical
behavior.

\subsection{Trust boundary, blockers, and negative controls}

MPFR, GMP, native build, Rust compiler, C ABI and FFI, \path{gmp-mpfr-sys}, Rug, wrapper code,
expression construction, and data identity remain trusted.  The \path{forbid(unsafe_code)}
attribute in \path{pid-core} does not verify native transitive dependencies.

Required negative controls must reverse one endpoint direction, retain an unresolved near-zero
atom, retain an unresolved minimizer order, report a precision limit, bind a count mutation to a
new digest, and compare small tables with an independent high-precision implementation.

\subsection{Permitted claim}

\begin{quote}
The certified backend returns an outward enclosure of the exact-real expression under MPFR
4.2.2's directed-rounding contract and the audited wrapper.  It does not prove the existing
binary64 result, the Rust compiler or FFI, or the MPFR implementation itself.
\end{quote}

Do not state that MPFR formally verified the result, that the current \path{f64} output is
certified, or that an interval proves estimator calibration.

\section{Rocq Interval adoption record}

\subsection{Upstream pin and license}

The selected stable pins are Rocq 9.2.0, released on 27 March 2026, and
\path{coq-interval} 4.11.4, published on 23 February 2026.  Rocq 9.3 release-candidate builds are
prereleases and are not selected.

Rocq core is LGPL-2.1-only.  \path{coq-interval} is CeCILL-C.

The package archive SHA-512 for \path{coq-interval} 4.11.4 is:

\begin{lstlisting}
5b2ccff32c3d6caa9002455dc472d6c4c8f70337581d038ca432dd92ed90b262
b203960a2a62896c12992be2ff2436bf7e6a8ffd3aec625859f47088d262ef00
\end{lstlisting}

The two lines are one continuous digest.

Primary sources are:
\begin{itemize}[leftmargin=*]
  \item \href{https://rocq-prover.org/}{Rocq site};
  \item \href{https://github.com/rocq-prover/rocq/releases/tag/V9.2.0}
  {Rocq 9.2.0 release};
  \item \href{https://raw.githubusercontent.com/rocq-prover/rocq/V9.2.0/rocq-core.opam}
  {core package metadata};
  \item \href{https://rocq-prover.org/p/coq-interval/4.11.4}{Interval package record};
  \item \href{https://rocq-prover.org/p/coq-interval/latest/versions}{Interval versions};
  \item \href{https://coqinterval.gitlabpages.inria.fr/}{Interval documentation};
  \item
  \href{https://raw.githubusercontent.com/coq/opam-coq-archive/master/released/packages/coq-interval/coq-interval.4.11.4/opam}
  {exact OPAM metadata};
  \item \href{https://rocq-prover.org/docs/using-opam}{Rocq OPAM instructions}; and
  \item
  \href{https://docs.rocq-prover.org/master/refman/practical-tools/coq-commands.html}
  {compile and check commands}.
\end{itemize}

The exact OPAM metadata permits the modern split \path{coq-core} and \path{coq-stdlib} route.  It
also requires Flocq, MathComp, Coquelicot, bignums, OCaml, and a compiler.

\subsection{Pin and installation rule}

\begin{lstlisting}
opam switch create 4.14.2
eval "$(opam env)"
opam repo add rocq-released https://rocq-prover.org/opam/released
opam install rocq-prover=9.2.0 rocq-core=9.2.0 coq-interval=4.11.4
opam pin add rocq-core 9.2.0
opam pin add coq-interval 4.11.4
opam switch export --full rocq-9.2.0-interval-4.11.4.export
\end{lstlisting}

Retain exact OPAM version, repository commit, full switch export, package hashes, certificate
source digest, compiled object digest, and checker output with assumptions.

\begin{lstlisting}
rocq compile -q Certificate.v
rocq check -o Certificate.vo
\end{lstlisting}

The \path{-o} output exposes the logical context and assumptions.  The evidence gate must reject
\path{Admitted}, \path{admit}, unexpected axioms, and an incomplete check.

\subsection{First obligations}

Use Rocq Interval only after the interval expression has a versioned schema.  Candidate theorems
can state:
\begin{itemize}[leftmargin=*]
  \item a rational logarithmic expression lies inside emitted endpoints;
  \item an atom interval is strictly positive or negative;
  \item two \(\mathrm{I}_{\min}\) intervals are strictly ordered;
  \item a M{\"o}bius reconstruction encloses each cumulative coordinate; and
  \item a certificate binds to exact count-table and expression-schema digests.
\end{itemize}

The theorem statement must bind to the exact input identity.  A theorem about the wrong count
table can still be a valid theorem.

\subsection{Trust boundary, blockers, and negative controls}

Rocq's kernel checks the proof object.  It does not prove that the generator encoded the intended
formula or data, that a file digest was bound to a runtime operation, that \texttt{pid-rs} used the
result, or that the statement has statistical meaning.

The generator, statement schema, input binding, serialization code, OCaml runtime, Rocq kernel,
and platform remain in the trust boundary.

Required negative controls must mutate an endpoint, mutate the bound input digest, insert an
admission, add an unexpected axiom, and present a valid proof for a different count table.  Each
control must fail the applicable proof or identity gate.

\subsection{Permitted claim}

\begin{quote}
Rocq's kernel independently checked the emitted inequality theorem and exposed its remaining
assumptions.  It does not prove that the generator encoded the intended data and expression or
that the \texttt{pid-rs} runtime used that certificate.
\end{quote}

Do not state that Rocq verified the Rust implementation, proved input provenance, or proved
application validity.

\section{Cross-layer evidence rules}

\subsection{A result stays inside its layer}

\small
\begin{longtable}{L{0.18\textwidth}L{0.30\textwidth}L{0.39\textwidth}}
\toprule
Evidence & What it can support & What it cannot support by itself \\
\midrule
\endfirsthead
\toprule
Evidence & What it can support & What it cannot support by itself \\
\midrule
\endhead
Existing Lean theorem & Exact theorem in the declared Lean model &
Rust implementation, binary64 output, or data assumptions \\
Verus proof & Contracts for opted-in supported functions &
Compiler, external assumptions, or transcendental accuracy \\
Kani harness & Exhaustive bounded property &
Unbounded inputs, concurrent behavior, or numerical accuracy \\
MPFR interval & Outward enclosure for one encoded expression &
Existing \path{f64}, statistical calibration, or an FFI proof \\
Rocq certificate & Kernel-checked encoded theorem &
Correct generator, runtime use, or intended data identity \\
Hidden benchmark & Observed behavior on a frozen corpus &
Universal correctness or a mathematical theorem \\
\bottomrule
\end{longtable}
\normalsize

Evidence can compose only when a checked bridge binds source revision, tool revision,
specification revision, input identity, generated-artifact identity, result identity, assumptions,
and claim scope.

\subsection{Required evidence envelope}

Every accepted pilot result must record:
\begin{enumerate}[leftmargin=*]
  \item exact tool version or commit;
  \item compiler, solver, and backend pins;
  \item host, target, and installation asset digest;
  \item source commit and dirty-state policy;
  \item command line;
  \item input bounds or theorem statement;
  \item assumption inventory;
  \item generated artifact hashes;
  \item checker result;
  \item negative-control result;
  \item unsupported cases; and
  \item exact permitted claim text.
\end{enumerate}

A green exit code without this envelope is a development result.  It is not release evidence.

\subsection{Fail-closed rules}

The process must return a typed failure or unresolved result if a pin differs, a digest is absent,
a solver or compiler version is unknown, an unwinding assertion fails, an assumption is not
approved, an artifact digest differs, an interval does not determine the requested result, a
certificate has an admission or unexpected axiom, or a supported-subset translator rejects the
source.

The process must not change a failure into an approximate success.

\section{Pilot exit criteria}

\paragraph{Rug and MPFR.}
All endpoints must be outward rounded.  Precision escalation must have a typed result.  Ambiguous
signs and ties must stay unresolved.  Exact identities and independent small-table checks must be
retained.  A license review must define the distribution route.

\paragraph{Kani.}
Every harness must declare bounds.  Every loop must have a passing unwinding assertion.  Meaningful
mutations must produce counterexamples.  No harness can model transcendental accuracy.  The report
must state the bounded domain.

\paragraph{Verus.}
The complete selected kernel must verify.  The assumption inventory must be empty or approved.
Semantic mutations must fail.  The erased program must pass Rust 1.89 gates.  Ordinary tests must
agree with specification fixtures.

\paragraph{Rocq Interval.}
The schema must be versioned.  \path{rocq check -o} must report no unexpected assumption.  The gate
must reject admissions.  Identity mutations must fail.  Archived exact inputs must regenerate the
certificate.

\paragraph{Aeneas.}
The selected subset must translate without an unsound workaround.  External models must be minimal
and reviewed.  The theorem correspondence must be documented.  Lean compatibility must be
resolved.  Source mutations must affect the model and proof.

\section{Conclusion}

The most useful immediate addition is an executable categorical interval oracle with explicit
unresolved states.  Kani can then attack bounded compiled behavior.  Verus can attack unbounded
functional contracts for the same small kernel.

Rocq Interval is useful only when an independent certificate has operational value and the
generator-to-statement bridge is controlled.  Aeneas is a valuable research route, but its current
release and toolchain state does not justify production adoption.

\begin{quote}
\texttt{pid-rs} is not end-to-end formally verified.  Each future result must state the exact
layer, theorem or bound, assumptions, tool pins, trust boundary, and unsupported claims.
\end{quote}

\section{Primary source index}

\subsection*{Verus}
\begin{itemize}[leftmargin=*]
  \item
  \href{https://github.com/verus-lang/verus/releases/tag/release/0.2026.07.18.3a4d30b}
  {Release \path{0.2026.07.18.3a4d30b}}
  \item \href{https://verus-lang.github.io/verus/guide/}{Verus guide}
  \item \href{https://verus-lang.github.io/verus/guide/cargo_verus.html}{Cargo integration}
  \item \href{https://verus-lang.github.io/verus/guide/tcb.html}{Trusted computing base}
  \item
  \href{https://verus-lang.github.io/verus/guide/reference-assume-specification.html}
  {Assumption specifications}
\end{itemize}

\subsection*{Kani}
\begin{itemize}[leftmargin=*]
  \item \href{https://github.com/model-checking/kani/releases/tag/kani-0.67.0}
  {Release 0.67.0}
  \item \href{https://model-checking.github.io/kani/install-guide.html}{Installation guide}
  \item \href{https://model-checking.github.io/kani/}{Kani manual}
  \item \href{https://model-checking.github.io/kani/tutorial-loop-unwinding.html}{Loop unwinding}
  \item \href{https://model-checking.github.io/kani/rust-feature-support.html}
  {Rust feature support}
  \item \href{https://model-checking.github.io/kani/rust-feature-support/intrinsics.html}
  {Intrinsic support}
\end{itemize}

\subsection*{Aeneas and Charon}
\begin{itemize}[leftmargin=*]
  \item \href{https://github.com/AeneasVerif/aeneas}{Aeneas repository}
  \item \href{https://github.com/AeneasVerif/aeneas/releases}{Aeneas releases}
  \item
  \href{https://github.com/AeneasVerif/aeneas/commit/6e167c9b63a4dafd66d0e0edd9d94669f957ff7b}
  {Pinned Aeneas commit}
\end{itemize}

\subsection*{MPFR and Rug}
\begin{itemize}[leftmargin=*]
  \item \href{https://www.mpfr.org/}{MPFR project}
  \item \href{https://www.mpfr.org/mpfr-current/}{MPFR 4.2.2}
  \item \href{https://www.mpfr.org/mpfr-current/mpfr.html}{MPFR manual}
  \item \href{https://docs.rs/rug/1.30.0/rug/}{Rug 1.30.0}
  \item \href{https://docs.rs/gmp-mpfr-sys/1.7.1/gmp_mpfr_sys/}
  {\path{gmp-mpfr-sys} 1.7.1}
  \item \href{https://gmplib.org/manual/Copying}{GMP copying conditions}
\end{itemize}

\subsection*{Rocq and Interval}
\begin{itemize}[leftmargin=*]
  \item \href{https://rocq-prover.org/}{Rocq project}
  \item \href{https://github.com/rocq-prover/rocq/releases/tag/V9.2.0}{Rocq 9.2.0}
  \item \href{https://rocq-prover.org/p/coq-interval/4.11.4}
  {\path{coq-interval} 4.11.4}
  \item \href{https://coqinterval.gitlabpages.inria.fr/}{Interval documentation}
  \item \href{https://rocq-prover.org/docs/using-opam}{Rocq OPAM instructions}
  \item
  \href{https://docs.rocq-prover.org/master/refman/practical-tools/coq-commands.html}
  {Rocq compile and check tools}
\end{itemize}

\end{document}
